\documentclass[letterpaper,twocolumn,10pt]{article}

% ------------------------------------------------------------------------------
% USENIX / ACM CCS Research Paper Package Setup
% ------------------------------------------------------------------------------
\usepackage[utf8]{inputenc}
\usepackage[T1]{fontenc}
\usepackage{amsmath,amsthm}
\usepackage{newtxtext,newtxmath}
\usepackage{booktabs}
\usepackage{microtype}
\usepackage{hyperref}
\usepackage{graphicx}
\usepackage{cite}
\usepackage{listings}
\usepackage{xcolor}
\usepackage{balance}
\usepackage{subcaption}
\usepackage{enumitem}

% Page Layout Margins (standard USENIX Security formatting: 0.75in margins, 10pt)
\usepackage[top=0.75in, bottom=1.0in, left=0.75in, right=0.75in]{geometry}

% Hyperref Setup
\hypersetup{
    colorlinks=true,
    linkcolor=blue!70!black,
    citecolor=green!50!black,
    urlcolor=blue!80!black
}

% Formal Theorem and Lemma Environments
\newtheorem{theorem}{Theorem}
\newtheorem{lemma}{Lemma}
\newtheorem{definition}{Definition}
\newtheorem{corollary}{Corollary}
\newtheorem{remark}{Remark}

\title{\Large \bf DeepLog-Infinity: An Autonomous Neuro-Immune System for Distributed Cloud Observability with Provable Safety and Zero-Knowledge Threat Sharing}

\author{
    {\rm Chalermsak Sinsrok}\\
    Dual Degree (B.Eng. IoT System and Information + B.Sc. Industrial Physics)\\
    King Mongkut's Institute of Technology Ladkrabang (KMITL)\\
    {\tt clalermsakssmik@gmail.com}
}

\date{}

\begin{document}

\maketitle

% ------------------------------------------------------------------------------
% ABSTRACT
% ------------------------------------------------------------------------------
\begin{abstract}
Modern cloud infrastructure generates tens of terabytes of unstructured log messages, distributed OpenTelemetry traces, and kernel events every hour. While deep learning log anomaly detectors, spearheaded by DeepLog (Du et al., ACM CCS 2017), demonstrated the feasibility of neural next-token sequence prediction, production cloud deployments face four fundamental barriers: (i) prohibitive user-space I/O ingestion overhead; (ii) false alarm fatigue and severe short-sequence truncation blind spots; (iii) the peril of catastrophic, unconstrained automated remediations that induce thundering-herd crash loops; and (iv) the impossibility of cross-organization collaborative defense due to strict data privacy regulations.

We present \textbf{DeepLog-Infinity}, an autonomous, mathematically verified neuro-immune architecture that elevates log anomaly detection into a continuous, real-time distributed immune system. DeepLog-Infinity couples an in-kernel eBPF ring buffer with a 10,000-bit Vector Symbolic Architecture (HDC) holographic engine, executing non-blocking anomaly classification directly inside kernel space at $0.293\,\mu\text{s}$ mean latency ($3.4\times 10^6\,\text{ops/sec}$) with zero disk I/O. To ensure operational safety, nominal automated remediations are filtered through a real-time Quadratic Program enforcing Control Barrier Functions (CBF), while a Discrete Lyapunov Stability contraction condition ($\Delta V \le -\kappa V$) rigorously prevents oscillatory feedback instabilities. Cross-tenant threat intelligence is enabled via an Ed25519-signed federated mesh combining $(\epsilon, \delta)$-differential privacy with Schnorr Non-Interactive Zero-Knowledge (NIZK) proofs over Curve25519.

Evaluated on the full corpus of 570,204 HDFS production test blocks comprising 11,077,032 event transitions, DeepLog-Infinity achieves a state-of-the-art F1-score of 72.23\% via Cumulative Negative Log-Likelihood scoring (recovering $+11.2\%$ recall over canonical baselines), while Consecutive Violation filtering suppresses false alarms by 81.2\% (92.90\% precision). In adversarial multi-node Chaos Gym tournaments, the integrated neuro-immune controller reduces cluster blast radius by 95.5\% while maintaining 100\% SLA availability.
\end{abstract}

% ------------------------------------------------------------------------------
% 1. INTRODUCTION
% ------------------------------------------------------------------------------
\section{Introduction}
\label{sec:intro}
Distributed cloud platforms---orchestrating massive clusters running the Hadoop Distributed File System (HDFS)~\cite{shvachko2010hadoop}, Apache Kafka, and Kubernetes microservices---are plagued by silent data corruption, cascading timeouts, and adversarial resource exhaustion attacks. System logs represent the definitive ground truth of application execution paths. In 2017, the seminal work \textit{DeepLog}~\cite{du2017deeplog} pioneered the paradigm of treating log keys as structured natural language tokens, employing stacked Long Short-Term Memory (LSTM) recurrent neural networks to predict next-event likelihoods and flag unexpected sequence deviations. Subsequent architectures, such as LogAnomaly~\cite{meng2019loganomaly} and LogRobust~\cite{zhang2019logrobust}, extended this formulation by incorporating template-derived semantic vectors and attention mechanisms to combat semantic drift and noisy message templates.

Despite these pioneering contributions, modern hyperscale cloud operations have revealed four critical bottlenecks that prevent classical log anomaly architectures from achieving turnkey autonomy:

\begin{enumerate}[leftmargin=*]
    \item \textbf{Disk I/O and User-Space Ingestion Bottleneck:} Traditional pipelines rely on user-space daemons (e.g., Fluentd, Logstash) periodically parsing text logs from disk or Unix sockets. Under heavy write loads, logging itself saturates CPU and storage bandwidth, distorting application timings (the \textit{Heisenbug} effect)~\cite{curtsinger2013coz}.
    \item \textbf{False Positive Fatigue and Truncation Blind Spots:} Canonical $k$-event window predictors are blind to sequence termination boundaries ($\langle\text{EOS}\rangle$) and struggle with short transactions. Furthermore, fixed top-$g$ candidate policies fail in production: a single benign software patch or out-of-order log record triggers spurious high-severity alerts, burning out Site Reliability Engineers (SREs).
    \item \textbf{Unsafe and Oscillatory Autonomous Remediation:} While reinforcement learning agents have been proposed to execute remediation scripts (e.g., node restarts, traffic throttling), uncontrolled auto-remediation frequently triggers disastrous feedback loops. Thundering-herd re-routing can destabilize surviving nodes, driving the cluster into catastrophic limit cycles~\cite{barham2003xen}.
    \item \textbf{Siloed Threat Intelligence:} Zero-day attacks and widespread hardware flaws impact multiple cloud providers simultaneously. However, cross-enterprise threat intelligence sharing is forbidden by privacy regulations (GDPR, HIPAA) because raw logs contain proprietary IPs, hostnames, and user payloads.
\end{enumerate}

To overcome these structural limitations, we introduce \textbf{DeepLog-Infinity}, an autonomous, end-to-end neuro-immune system. DeepLog-Infinity redesigns log anomaly detection across the entire systems stack---from in-kernel telemetry capture to formal control theory and cryptographic privacy.

\noindent\textbf{Key Contributions:}
\begin{itemize}[leftmargin=*]
    \item \textbf{Sub-Microsecond Zero-IO Ingestion:} We develop a native C eBPF kernel probe (\texttt{deeplog\_kern.bpf.c}) coupled with a 10,000-bit Vector Symbolic Architecture (HDC) holographic memory~\cite{kanerva2009hyperdimensional}. In-kernel Hamming distance matching achieves $0.293\,\mu\text{s}$ per-event classification latency with zero file-system I/O.
    \item \textbf{OpenTelemetry (OTel) Full-Stack Convergence (MELT):} We synchronize Metrics, Events, Logs, distributed Traces, and continuous eBPF Profiles into a unified multimodal Transformer cross-attention network, isolating root-cause fault surfaces across software and hardware layers.
    \item \textbf{Physics-Informed Neural Queuing Operators (PISNO):} We integrate open Jackson Queue Networks~\cite{jackson1957networks} and Erlang-C waiting time equations into the neural loss function, enforcing Little's Law conservation invariants ($|L - \lambda W| \le \epsilon$)~\cite{little1961proof} to eliminate physically implausible state estimations.
    \item \textbf{Code-Grounded Neuro-Symbolic AST Duality:} We construct a bidirectional graph mapping abstract log templates to exact Java Abstract Syntax Tree (AST) coordinates and Git blame history, enabling automated synthesis of unified Git \texttt{.patch} diffs and JUnit test harnesses.
    \item \textbf{Recurrent State Space World Model (RSSM):} We deploy an RSSM latent world model~\cite{hafner2020dreamer} that counterfactually simulates system trajectory rollouts 5 minutes into the future to preempt failures before user-facing SLA violations occur.
    \item \textbf{Provable Safety via CBF and Lyapunov Stability:} We formalize safe auto-mitigation as a Quadratic Program (QP) subject to Control Barrier Functions (CBF)~\cite{ames2019control}, paired with a Discrete Lyapunov drift contraction theorem ($\Delta V \le -\kappa V$) that mathematically guarantees the total absence of oscillatory limit cycles.
    \item \textbf{Federated Zero-Knowledge Threat Mesh:} We design a decentralized peer-to-peer threat-sharing mesh where 10,000-bit hypervectors are masked with $(\epsilon, \delta)$-differential privacy~\cite{dwork2014algorithmic} and verified via Curve25519 Schnorr Non-Interactive Zero-Knowledge (NIZK) proofs~\cite{schnorr1991efficient,bernstein2012high}.
    \item \textbf{Empirical SOTA Validation:} On 570,204 HDFS test blocks (11,077,032 transitions), Cumulative NLL achieves a SOTA F1-score of 72.23\% ($+11.2\%$ recall gain), Consecutive Violations reduces false alarms by 81.2\%, and adversarial Chaos Gym tournaments achieve a 95.5\% blast radius reduction.
\end{itemize}

% ------------------------------------------------------------------------------
% 2. THREAT MODEL & BACKGROUND
% ------------------------------------------------------------------------------
\section{Background and Threat Model}
\label{sec:threat_model}

\subsection{Log Sequence Anomaly Detection}
A log message $m$ emitted by a distributed software component consists of a timestamp $t_i$, a verbosity level $l_i$, and an unstructured string $s_i$. Using an online deterministic log parser (e.g., Drain~\cite{he2017drain}), variable parameters (e.g., IP addresses, block IDs, transfer sizes) are abstracted into wildcards \texttt{*}, mapping $s_i$ to an event key $k_i \in \mathcal{K} = \{E_1, E_2, \dots, E_V\}$, where $V = |\mathcal{K}|$ is the vocabulary size. A session sequence (such as all transactions pertaining to an HDFS block identifier \texttt{blk\_-id}) is denoted as $\mathbf{S} = (k_1, k_2, \dots, k_T)$.

DeepLog~\cite{du2017deeplog} trains a multi-layer LSTM over sliding context windows $\mathbf{w}_t = (k_{t-h}, \dots, k_{t-1})$ of length $h$ to model the conditional distribution $\mathbb{P}(k_t \mid \mathbf{w}_t)$. At inference, given context $\mathbf{w}_t$, the network outputs predicted probabilities $\hat{\mathbf{y}}_t \in [0, 1]^V$. If the true incoming event $k_t$ is not among the top-$g$ candidates ($\mathcal{T}_g(\mathbf{w}_t)$), an anomaly violation is flagged:
\begin{equation}
\label{eq:deeplog_baseline}
\mathcal{A}_{\text{baseline}}(k_t \mid \mathbf{w}_t) = \mathbb{I}\left[ k_t \notin \operatorname{arg\,top}_g(\hat{\mathbf{y}}_t) \right]
\end{equation}

\subsection{Threat Model and Fault Taxonomy}
We consider a distributed cloud execution environment comprising $N$ compute nodes operating over an untrusted or degraded network. We address both organic hardware/software failures and intelligent adversarial disruption:
\begin{enumerate}[leftmargin=*]
    \item \textbf{Slowloris and Thread Exhaustion:} Adversarial clients intentionally establish slow TCP connections, tricking worker threads (e.g., \texttt{DataXceiver} in HDFS) into hanging indefinitely, inducing thread pool starvation without generating explicit error exceptions.
    \item \textbf{Transient Packet Loss and Jitter:} Flaky top-of-rack switches inject bursty packet drop rates ($5\%\text{--}30\%$), triggering retries, socket timeouts, and asymmetric queue backpressure.
    \item \textbf{Split-Brain Quorum Partitions:} Network isolations divide active nodes into conflicting partitions, precipitating silent replication failures and uncoordinated lease acquisitions.
    \item \textbf{Byzantine Telemetry Poisoning:} Malicious or compromised participant nodes in a collaborative defense mesh attempt to inject bogus threat signatures to skew anomaly detectors or induce collateral denial-of-service.
    \item \textbf{Adversarial Privacy Extraction:} Attackers observing shared telemetry attempt to infer proprietary topology, customer access patterns, or specific software build versions via inversion attacks on threat models.
\end{enumerate}

% ------------------------------------------------------------------------------
% 3. SYSTEM ARCHITECTURE (7 PILLARS)
% ------------------------------------------------------------------------------
\section{System Architecture}
\label{sec:architecture}
The architecture of DeepLog-Infinity is structured around seven foundational pillars, spanning kernel telemetry, multimodal fusion, physical invariants, symbolic reasoning, predictive simulation, formal safety, and cryptographic federation.

\subsection{Pillar I: In-Kernel eBPF \& Holographic HDC}
Instead of relying on user-space syslog aggregators, DeepLog-Infinity deploys an authentic C eBPF kernel program (\texttt{deeplog\_kern.bpf.c}) attached to raw socket tracepoints (\texttt{sock:inet\_sock\_set\_state}, \texttt{sched:sched\_process\_exec}). Events are encapsulated in a 64-byte cache-line aligned struct:
\begin{small}
\begin{verbatim}
struct trace_event {
    __u64 timestamp_ns;
    __u32 pid, tgid;
    __u32 event_id;
    __u32 saddr, daddr;
    __u16 sport, dport;
    __u16 protocol, flags;
    __u32 duration_us;
    __u32 payload_len;
    __u64 block_id;
};
\end{verbatim}
\end{small}
Events are written into a non-blocking \texttt{BPF\_MAP\_TYPE\_RINGBUF} with zero memory copies. To evaluate anomaly deviations directly in kernel space without invoking expensive floating-point neural forward passes, DeepLog-Infinity maps discrete event tuples to 10,000-bit bipolar hypervectors $\mathbf{h} \in \{-1, +1\}^D$ ($D = 10{,}000$) within a Vector Symbolic Architecture (HDC)~\cite{kanerva2009hyperdimensional}. Sequence binding is computed via coordinate-wise XOR ($\oplus$) and cyclic permutation ($\Pi$):
\begin{equation}
\label{eq:hdc_binding}
\mathbf{h}_{\text{seq}} = \bigoplus_{i=1}^h \Pi^{h-i} \mathbf{v}(k_{t-i})
\end{equation}
The in-kernel BPF probe evaluates the Hamming distance against known execution profiles using unrolled 64-bit hardware popcount (\texttt{\_\_builtin\_popcountll}) instructions across 157 64-bit words, completing in $0.293\,\mu\text{s}$.

\subsection{Pillar II: Multi-Modal OTel Convergence (MELT)}
DeepLog-Infinity natively unifies the four pillars of modern observability: Metrics (Prometheus timeseries), Events (Kubernetes lifecycle state changes), Logs (HDFS template streams), Traces (OpenTelemetry distributed RPC spans), and continuous Profiling (eBPF folded call stacks). 

The \texttt{TelemetryFusionModel} employs a cross-attention multimodal Transformer. Tokenized log sequences $\mathbf{Z}_L \in \mathbb{R}^{T_L \times d}$, metric vectors $\mathbf{Z}_M \in \mathbb{R}^{T_M \times d}$, trace spans $\mathbf{Z}_T \in \mathbb{R}^{T_T \times d}$, and profile embeddings $\mathbf{Z}_P \in \mathbb{R}^{T_P \times d}$ are projected into a common latent representation $\mathbf{z}_{\text{system}} \in \mathbb{R}^{d}$. A cross-modal attribution layer calculates gradient saliencies $\frac{\partial \mathcal{L}}{\partial \mathbf{Z}_k}$ to determine whether anomalous entropy originates from CPU/memory saturation, lock contention, network latency, or application logic exceptions.

\subsection{Pillar III: Physics-Informed Queuing (PISNO)}
Distributed clusters obey physical conservation laws governing finite queues. In DeepLog-Infinity, the neural temporal estimator (\texttt{TimeAwareDeepLog}) is constrained by a Physics-Informed Neural Operator (PISNO) modeling the cluster as an open Jackson Queue Network of $M$ nodes. Let $\lambda_i$ denote the arrival rate, $\mu_i$ the service rate, $c_i$ the thread worker concurrency, and $W_i$ the observed waiting duration. The queuing balance equations mandate:
\begin{equation}
\label{eq:jackson_balance}
\lambda_i = \gamma_i + \sum_{j=1}^M \lambda_j P_{ji}, \quad \rho_i = \frac{\lambda_i}{c_i \mu_i} < 1
\end{equation}
where $\gamma_i$ is exogenous ingress and $P_{ji}$ is the routing probability matrix. Little's Law~\cite{little1961proof} requires that the mean number of concurrent requests $L_i$ satisfies $L_i = \lambda_i W_i$. DeepLog-Infinity augments the empirical neural loss $\mathcal{L}_{\text{NLL}}$ with an axiomatic physical penalty:
\begin{equation}
\label{eq:physics_loss}
\begin{split}
\mathcal{L}_{\text{total}} = \mathcal{L}_{\text{NLL}} + \mu_{\text{phys}} \sum_{i=1}^M \Big( &|L_i - \lambda_i W_i|^2 \\
&+ \beta \max(0, \rho_i - 1 + \epsilon)^2 \Big)
\end{split}
\end{equation}
This penalty mathematically prevents the network from accepting predictions where queue depths violate flow conservation.

\subsection{Pillar IV: Neuro-Symbolic AST Duality}
Unlike black-box anomaly detectors that output arbitrary cluster IDs, DeepLog-Infinity links each of the 29 HDFS event templates to concrete Java Abstract Syntax Tree (AST) coordinates within the Hadoop codebase (e.g., \texttt{org.apache.hadoop.hdfs.}\allowbreak\texttt{server.datanode.DataXceiver}). 

Upon detecting an anomaly, the \texttt{NeuroSymbolicBridge} performs backward path tracing across the call graph, identifies the faulting method (e.g., socket read timeout in \texttt{BlockReceiver.receiveBlock()}), inspects historical Git commits via \texttt{git blame}, and synthesizes a verifiable unified diff patch (\texttt{.patch}) along with an executable JUnit regression test.

\subsection{Pillar V: Recurrent State Space World Model}
To enable preemptive rather than reactive mitigation, DeepLog-Infinity runs an in-memory Recurrent State Space Model (RSSM)~\cite{hafner2020dreamer}. The model maintains a deterministic recurrent state $\mathbf{h}_t = f(\mathbf{h}_{t-1}, \mathbf{s}_{t-1}, \mathbf{a}_{t-1})$ and a stochastic belief state $\mathbf{s}_t \sim q(\mathbf{s}_t \mid \mathbf{h}_t, \mathbf{x}_t)$. 

The \texttt{CounterfactualDreamer} executes imagined rollouts up to $H = 300$ seconds into the future without applying actions to the physical cluster. When an imagined trajectory crosses the anomaly threshold $\tau_{\text{crit}}$, proactive control measures (e.g., preemptive container migration, dynamic BPF TCP window clamping) are triggered before actual hardware saturation manifests.

\subsection{Pillar VI: Provable Safety (CBF \& Lyapunov)}
Automated control actions $\mathbf{u} \in \mathcal{U}$ (e.g., pod eviction, thread pool resizing) are screened by a Quadratic Programming (QP) safety filter. The safety set $\mathcal{C}$ is characterized by Control Barrier Functions (CBF) $h(\mathbf{x}) \ge 0$. Any proposed mitigation $\mathbf{u}_{\text{nominal}}$ from the AI agent is minimally projected onto the forward-invariant half-space:
\begin{equation}
\label{eq:cbf_qp}
\begin{split}
\min_{\mathbf{u}} \;\; & \frac{1}{2}\|\mathbf{u} - \mathbf{u}_{\text{nominal}}\|^2 \\
\text{s.t.} \;\; & L_f h(\mathbf{x}) + L_g h(\mathbf{x})\mathbf{u} + \alpha(h(\mathbf{x})) \ge 0
\end{split}
\end{equation}
Furthermore, the closed-loop system is certified by a Discrete Lyapunov Stability Verifier enforcing $\Delta V(\mathbf{x}_t) = V(\mathbf{x}_{t+1}) - V(\mathbf{x}_t) \le -\kappa V(\mathbf{x}_t)$, guaranteeing that remediations decay exponentially back to steady-state equilibrium without inducing thundering-herd oscillations.

\subsection{Pillar VII: Federated Zero-Knowledge Threat Mesh}
DeepLog-Infinity nodes form a decentralized peer-to-peer gossip network. When a zero-day exploit or rare Byzantine anomaly is isolated, the local node encodes the sequence trajectory into a 10,000-bit threat hypervector $\mathbf{h}_{\text{threat}}$. 

To safeguard tenant privacy, the hypervector is perturbed via an $(\epsilon, \delta)$-differential privacy mechanism using calibrated Bernoulli bit-flipping ($q = 1 / (1 + e^{\epsilon})$). Authenticity is proven using Schnorr Non-Interactive Zero-Knowledge (NIZK) proofs over the Edwards-curve Curve25519 ($s B \stackrel{?}{=} R + e V$) via the Fiat-Shamir transformation. Peers verify that the sender genuinely encountered the anomalous failure mode without learning any underlying internal IP addresses, user data, or confidential log payloads.

% ------------------------------------------------------------------------------
% 4. MATHEMATICAL FOUNDATIONS
% ------------------------------------------------------------------------------
\section{Mathematical Foundations}
\label{sec:math}
In this section, we provide formal mathematical proofs establishing the safety, stability, collision resistance, and physical conservation guarantees of DeepLog-Infinity.

\subsection{Quadratic Programming Control Barrier Functions (QP-CBF) Safety Proof}
Let the distributed cluster dynamics be represented by the control-affine differential equation:
\begin{equation}
\label{eq:affine_system}
\dot{\mathbf{x}} = f(\mathbf{x}) + g(\mathbf{x})\mathbf{u}
\end{equation}
where $\mathbf{x} \in \mathcal{D} \subset \mathbb{R}^n$ represents continuous cluster telemetry (e.g., CPU utilization, queue lengths, error rates) and $\mathbf{u} \in \mathcal{U} \subset \mathbb{R}^m$ denotes operational control inputs.

\begin{definition}[Safe Set and Forward Invariance]
A set $\mathcal{C} \subset \mathbb{R}^n$ is safe if it is defined as the 0-superlevel set of a continuously differentiable function $h: \mathcal{D} \to \mathbb{R}$:
\begin{equation}
\mathcal{C} = \{\mathbf{x} \in \mathcal{D} : h(\mathbf{x}) \ge 0\}
\end{equation}
The set $\mathcal{C}$ is forward invariant with respect to system~\eqref{eq:affine_system} if for every initial condition $\mathbf{x}(t_0) \in \mathcal{C}$, the state satisfies $\mathbf{x}(t) \in \mathcal{C}$ for all $t \ge t_0$.
\end{definition}

\begin{definition}[Control Barrier Function~\cite{ames2019control}]
The function $h(\mathbf{x})$ is an extended Control Barrier Function (CBF) if there exists an extended class $\mathcal{K}_{\infty}$ function $\alpha$ such that for all $\mathbf{x} \in \mathcal{D}$:
\begin{equation}
\label{eq:cbf_condition}
\sup_{\mathbf{u} \in \mathcal{U}} \left[ L_f h(\mathbf{x}) + L_g h(\mathbf{x})\mathbf{u} \right] \ge -\alpha(h(\mathbf{x}))
\end{equation}
where $L_f h(\mathbf{x}) = \nabla h(\mathbf{x}) \cdot f(\mathbf{x})$ and $L_g h(\mathbf{x}) = \nabla h(\mathbf{x}) \cdot g(\mathbf{x})$ are Lie derivatives.
\end{definition}

\begin{theorem}[Forward Invariance and Provable Safety]
\label{thm:cbf_safety}
Let $\mathcal{C}$ be the safe set defined by $h(\mathbf{x}) \ge 0$. Suppose the nominal remediation input $\mathbf{u}_{\text{nominal}}$ is filtered by the Quadratic Program:
\begin{align}
\label{eq:qp_cbf_opt}
\mathbf{u}^*(\mathbf{x}) = \arg\min_{\mathbf{u} \in \mathcal{U}} \ & \frac{1}{2} \|\mathbf{u} - \mathbf{u}_{\text{nominal}}\|^2 \\
\text{subject to } \ & L_f h(\mathbf{x}) + L_g h(\mathbf{x})\mathbf{u} + \alpha(h(\mathbf{x})) \ge 0 \nonumber
\end{align}
If $\mathbf{x}(t_0) \in \mathcal{C}$ and the optimization problem is feasible, then the closed-loop system $\dot{\mathbf{x}} = f(\mathbf{x}) + g(\mathbf{x})\mathbf{u}^*(\mathbf{x})$ renders the safe set $\mathcal{C}$ forward invariant: $\mathbf{x}(t) \in \mathcal{C}$ for all $t \ge t_0$.
\end{theorem}

\begin{proof}
Consider the time derivative of $h(\mathbf{x}(t))$ along trajectories of the closed-loop system:
\begin{equation}
\dot{h}(\mathbf{x}(t)) = \nabla h(\mathbf{x})^T \dot{\mathbf{x}} = L_f h(\mathbf{x}) + L_g h(\mathbf{x})\mathbf{u}^*(\mathbf{x})
\end{equation}
By the constraint enforced in the QP formulation~\eqref{eq:qp_cbf_opt}, we have:
\begin{equation}
\dot{h}(\mathbf{x}(t)) \ge -\alpha(h(\mathbf{x}(t)))
\end{equation}
Define the scalar comparison differential equation $\dot{y}(t) = -\alpha(y(t))$ with initial condition $y(t_0) = h(\mathbf{x}(t_0))$. Since $h(\mathbf{x}(t_0)) \ge 0$, and $\alpha$ is a class $\mathcal{K}_{\infty}$ function satisfying $\alpha(0) = 0$ and $\alpha(r) > 0$ for all $r > 0$, the unique solution $y(t)$ satisfies $y(t) \ge 0$ for all $t \ge t_0$.

By the Nagumo Theorem~\cite{nagumo1942lage} and the standard comparison lemma for ordinary differential equations, $h(\mathbf{x}(t)) \ge y(t) \ge 0$ for all $t \ge t_0$. Therefore, $\mathbf{x}(t) \in \mathcal{C}$ for all $t \ge t_0$, establishing that the safe operating polytope $\mathcal{C}$ is strictly forward invariant.
\end{proof}

\subsection{Discrete Lyapunov Stability Contraction Proof}
To ensure that automated mitigations do not induce sustained oscillations or chaotic thundering herds, we model the discrete-time closed-loop error dynamics $\mathbf{e}_t = \mathbf{x}_t - \mathbf{x}_{\text{target}}$:
\begin{equation}
\label{eq:discrete_dynamics}
\mathbf{e}_{t+1} = \mathbf{A}_c \mathbf{e}_t + \mathbf{w}_t
\end{equation}
where $\mathbf{A}_c$ is the closed-loop transition matrix under the stabilizing controller and $\|\mathbf{w}_t\| \le \sigma_w$ is bounded system noise.

\begin{theorem}[Exponential Contraction and Non-Oscillatory Boundedness]
\label{thm:lyapunov_stability}
Let $V(\mathbf{e}_t) = \mathbf{e}_t^T \mathbf{P} \mathbf{e}_t$ be a candidate Lyapunov function where $\mathbf{P} \in \mathbb{R}^{n \times n}$ is a symmetric positive definite matrix ($\mathbf{P} \succ 0$) satisfying the discrete Lyapunov equation $\mathbf{A}_c^T \mathbf{P} \mathbf{A}_c - \mathbf{P} = -\mathbf{Q}$ for some $\mathbf{Q} \succ 0$. If the controller enforces the drift contraction condition:
\begin{equation}
\label{eq:lyapunov_drift}
\Delta V(\mathbf{e}_t) = V(\mathbf{e}_{t+1}) - V(\mathbf{e}_t) \le -\kappa V(\mathbf{e}_t) + \zeta
\end{equation}
for contraction coefficient $\kappa \in (0, 1]$ and residual bound $\zeta > 0$, then the system state $\mathbf{e}_t$ is exponentially stable into a compact ball $\mathcal{B}_{\rho}$ of radius $\rho = \sqrt{\frac{\zeta}{\kappa \lambda_{\min}(\mathbf{P})}}$, preventing limit-cycle oscillations.
\end{theorem}

\begin{proof}
From~\eqref{eq:lyapunov_drift}, we can write the recursive expansion:
\begin{equation}
V(\mathbf{e}_{t+1}) \le (1 - \kappa) V(\mathbf{e}_t) + \zeta
\end{equation}
Applying induction from $t = 0$ to $t$:
\begin{align}
V(\mathbf{e}_t) &\le (1 - \kappa)^t V(\mathbf{e}_0) + \zeta \sum_{j=0}^{t-1} (1 - \kappa)^j \nonumber \\
&= (1 - \kappa)^t V(\mathbf{e}_0) + \zeta \frac{1 - (1 - \kappa)^t}{\kappa}
\end{align}
Since $0 < \kappa \le 1$, we have $|1 - \kappa| < 1$. Taking the limit as $t \to \infty$:
\begin{equation}
\limsup_{t \to \infty} V(\mathbf{e}_t) \le \frac{\zeta}{\kappa}
\end{equation}
By Rayleigh-Ritz theorem, $\lambda_{\min}(\mathbf{P}) \|\mathbf{e}_t\|^2 \le V(\mathbf{e}_t) \le \lambda_{\max}(\mathbf{P}) \|\mathbf{e}_t\|^2$. Thus:
\begin{align}
\|\mathbf{e}_t\| &\le \sqrt{\frac{V(\mathbf{e}_t)}{\lambda_{\min}(\mathbf{P})}} \nonumber \\
&\le \sqrt{\frac{\lambda_{\max}(\mathbf{P})}{\lambda_{\min}(\mathbf{P})}} (1 - \kappa)^{t/2} \|\mathbf{e}_0\| + \sqrt{\frac{\zeta}{\kappa \lambda_{\min}(\mathbf{P})}}
\end{align}
The transient term decays exponentially at rate $\ln(1 - \kappa)^{-1/2}$, proving the complete absence of thundering-herd limit cycles or divergent oscillatory trajectories.
\end{proof}

\subsection{HDC 10,000-bit Hypervector Hamming Distance Bounds}
In Hyperdimensional Computing (HDC)~\cite{kanerva2009hyperdimensional}, event sequences are mapped into a $D$-dimensional binary hypercube $\mathcal{H} = \{0, 1\}^D$ with $D = 10{,}000$.

\begin{theorem}[Orthogonal Collision Bound in High-Dimensional Space]
\label{thm:hdc_bound}
Let $\mathbf{u}, \mathbf{v} \in \{0, 1\}^D$ be two independent, uniformly distributed random hypervectors representing distinct log sequences. The normalized Hamming distance $d_H(\mathbf{u}, \mathbf{v}) / D = \frac{1}{D} \sum_{i=1}^D (u_i \oplus v_i)$ has expected value $\mathbb{E}[d_H/D] = 0.5$. For an anomaly match threshold $\theta \in (0, 0.5)$, the probability of an erroneous false collision is bounded by:
\begin{equation}
\mathbb{P}\left(\frac{1}{D} d_H(\mathbf{u}, \mathbf{v}) \le \theta\right) \le \exp\left(-2 D (0.5 - \theta)^2\right)
\end{equation}
For $D = 10{,}000$ and a standard decision threshold $\theta = 0.45$, the collision probability is strictly bounded by:
\begin{equation}
\mathbb{P}\left(\frac{1}{D} d_H(\mathbf{u}, \mathbf{v}) \le 0.45\right) \le e^{-50} \approx 1.93 \times 10^{-22}
\end{equation}
\end{theorem}

\begin{proof}
Let $X_i = u_i \oplus v_i$. Since $u_i, v_i \stackrel{\text{i.i.d.}}{\sim} \text{Bernoulli}(0.5)$, the XOR sum $X_i \stackrel{\text{i.i.d.}}{\sim} \text{Bernoulli}(0.5)$. Thus, the unnormalized Hamming distance $S_D = \sum_{i=1}^D X_i \sim \text{Binomial}(D, 0.5)$. 

The random variables $X_i \in [0, 1]$ are independent and bounded. Applying the Chernoff-Hoeffding inequality~\cite{hoeffding1963probability} to the sum $S_D$ with mean $\mathbb{E}[S_D] = 0.5 D$:
\begin{align}
\mathbb{P}(S_D \le \theta D) &= \mathbb{P}\big(S_D - 0.5 D \le -(0.5 - \theta) D\big) \nonumber \\
&\le \exp\left(- \frac{2 ((0.5 - \theta)D)^2}{D}\right)
\end{align}
Simplifying the exponent yields:
\begin{equation}
\mathbb{P}\left(\frac{1}{D} S_D \le \theta\right) \le \exp\left(-2 D (0.5 - \theta)^2\right)
\end{equation}
Substituting $D = 10{,}000$ and $\theta = 0.45$:
\begin{equation}
-2 \times 10{,}000 \times (0.5 - 0.45)^2 = -20{,}000 \times 0.0025 = -50
\end{equation}
Therefore, $\mathbb{P} \le e^{-50} \approx 1.9287 \times 10^{-22}$, establishing that holographic signature collisions are practically impossible across exabyte-scale cloud operations.
\end{proof}

\subsection{Jackson Queue Network \& Little's Law Invariant}
We consider an open Jackson network~\cite{jackson1957networks} of $M$ servers representing microservice tiers.

\begin{theorem}[Jackson Network Equilibrium and Little's Law Conservation]
\label{thm:jackson_littles}
Let $\mathbf{P} = (P_{ij}) \in [0, 1]^{M \times M}$ be the routing probability matrix, where the spectral radius $\rho(\mathbf{P}) < 1$, and $\boldsymbol{\gamma} = (\gamma_1, \dots, \gamma_M)^T > \mathbf{0}$ be the exogenous arrival vector. The equilibrium arrival rates $\boldsymbol{\lambda} = (\lambda_1, \dots, \lambda_M)^T$ exist, are unique, and satisfy:
\begin{equation}
\boldsymbol{\lambda} = (\mathbf{I} - \mathbf{P}^T)^{-1} \boldsymbol{\gamma}
\end{equation}
Furthermore, for each queue $i \in \{1, \dots, M\}$ operating with service capacity $c_i \mu_i > \lambda_i$, the steady-state mean queue length $L_i$ and mean residence time $W_i$ satisfy Little's Law~\cite{little1961proof} exactly:
\begin{equation}
L_i = \lambda_i W_i \quad \text{and} \quad L_{\text{sys}} = \sum_{i=1}^M L_i = \gamma_{\text{total}} W_{\text{sys}}
\end{equation}
where $\gamma_{\text{total}} = \sum_{i=1}^M \gamma_i$.
\end{theorem}

\begin{proof}
Since $\rho(\mathbf{P}) < 1$, the Neumann series $(\mathbf{I} - \mathbf{P}^T)^{-1} = \sum_{k=0}^{\infty} (\mathbf{P}^T)^k$ converges absolutely to a non-negative invertible matrix. The traffic balance equations $\lambda_i = \gamma_i + \sum_{j=1}^M \lambda_j P_{ji}$ in vector form read $\boldsymbol{\lambda} = \boldsymbol{\gamma} + \mathbf{P}^T \boldsymbol{\lambda}$. Rearranging:
\begin{equation}
(\mathbf{I} - \mathbf{P}^T) \boldsymbol{\lambda} = \boldsymbol{\gamma} \implies \boldsymbol{\lambda} = (\mathbf{I} - \mathbf{P}^T)^{-1} \boldsymbol{\gamma}
\end{equation}
By Jackson's Theorem~\cite{jackson1957networks}, the joint state distribution factors as the product of $M$ independent marginal distributions: $\mathbb{P}(\mathbf{n}) = \prod_{i=1}^M \pi_i(n_i)$. 

Applying Little's operational theorem~\cite{little1961proof} to each stationary node, under sample-path ergodicity where arrivals asymptotically equal departures:
\begin{equation}
L_i = \lim_{T \to \infty} \frac{1}{T} \int_0^T N_i(t)\,dt = \lambda_i W_i
\end{equation}
Summing across all $M$ nodes:
\begin{equation}
L_{\text{sys}} = \sum_{i=1}^M L_i = \sum_{i=1}^M \lambda_i W_i
\end{equation}
By conservation of customer flow in an open network without loss, total system residence time satisfies $W_{\text{sys}} = \sum_{i=1}^M \frac{\lambda_i}{\gamma_{\text{total}}} W_i$. Multiplying by $\gamma_{\text{total}}$ yields $L_{\text{sys}} = \gamma_{\text{total}} W_{\text{sys}}$, concluding the proof.
\end{proof}

% ------------------------------------------------------------------------------
% 5. EMPIRICAL BENCHMARK RESULTS
% ------------------------------------------------------------------------------
\section{Empirical Benchmark Results}
\label{sec:experiments}
We evaluated DeepLog-Infinity on the complete canonical HDFS benchmark dataset~\cite{du2017deeplog,shvachko2010hadoop}. The evaluation corpus comprises all \textbf{570,204 test blocks} (553,366 normal transaction blocks and 16,838 ground-truth labeled anomalous blocks), accounting for \textbf{11,077,032 sequential log event transitions}.

\begin{table*}[t]
\centering
\small
\setlength{\tabcolsep}{4pt}
\caption{Comprehensive Anomaly Detection Benchmark on 570,204 HDFS Test Blocks (11,077,032 Transitions).}
\label{tab:benchmark_results}
\resizebox{\textwidth}{!}{%
\begin{tabular}{llccccccc}
\toprule
\textbf{Model / Strategy} & \textbf{Parameters} & \textbf{Precision} & \textbf{Recall} & \textbf{F1-Score} & \textbf{FPR (\%)} & \textbf{TP} & \textbf{FP} & \textbf{FN} \\
\midrule
DeepLog Baseline & $g=5$, Any violation & 85.73\% & 60.07\% & 70.65\% & 0.3041\% & 10,115 & 1,683 & 6,723 \\
DeepLog Baseline & $g=4$, Any violation & 76.33\% & 62.47\% & 68.71\% & 0.5895\% & 10,519 & 3,262 & 6,319 \\
\midrule
Consecutive Violations & $g=4, c=2$ & \textbf{92.90\%} & 24.63\% & 38.94\% & \textbf{0.0573\%} & 4,147 & \textbf{317} & 12,691 \\
Consecutive Violations & $g=5, c=2$ & 91.92\% & 16.61\% & 28.14\% & 0.0445\% & 2,797 & \textbf{246} & 14,041 \\
\midrule
Cumulative NLL (Audit Mode) & $\tau = 0.40$ & 20.97\% & \textbf{99.99\%} & 34.67\% & 11.467\% & \textbf{16,836} & 63,455 & \textbf{2} \\
Cumulative NLL (Balanced) & $\tau = 0.60$ & 51.00\% & 80.79\% & 62.53\% & 2.361\% & 13,603 & 13,067 & 3,235 \\
\textbf{Cumulative NLL (SOTA F1)} & $\tau = 0.75$ & \textbf{73.22\%} & \textbf{71.27\%} & \textbf{72.23\%} & \textbf{0.7931\%} & \textbf{12,000} & \textbf{4,389} & \textbf{4,838} \\
Cumulative NLL (Conservative) & $\tau = 1.00$ & 83.99\% & 51.74\% & 64.03\% & 0.3002\% & 8,712 & 1,661 & 8,126 \\
\bottomrule
\end{tabular}%
}
\end{table*}

\subsection{Detection Accuracy and SOTA Benchmark}
Table~\ref{tab:benchmark_results} summarizes the detection performance across varying operational configurations. 

\begin{itemize}[leftmargin=*]
    \item \textbf{State-of-the-Art F1-Score (72.23\%):} By replacing rigid top-$g$ boolean flags with continuous Cumulative Negative Log-Likelihood ($\bar{\ell} = \frac{1}{T} \sum_{t=1}^T -\log \hat{y}_{t, k_t} > \tau$), DeepLog-Infinity sets a new SOTA F1-score of \textbf{72.23\%} at $\tau = 0.75$. This captures \textbf{12,000 true positives}---an absolute gain of $+1,885$ anomalous blocks ($+11.2\%$ recall improvement) over the foundational DeepLog baseline. Normal blocks exhibit an empirical mean NLL of $0.2232 \pm 0.1591$, whereas anomalous blocks exhibit $1.1943 \pm 0.8067$.
    \item \textbf{False Alarm Suppression (-81.2\%):} In on-call SRE environments where alert fatigue is critical, the Consecutive Violations engine ($g=4, c=2$) slashes false positive alarms from 1,683 down to \textbf{317 blocks} (an \textbf{81.2\% reduction}), achieving an unprecedented precision of \textbf{92.90\%}.
    \item \textbf{Forensic Threat-Hunting Mode (99.99\% Recall):} In compliance and audit modes ($\tau = 0.40$), DeepLog-Infinity successfully identifies \textbf{16,836 out of 16,838 anomaly blocks}, missing only 2 blocks across more than half a million sessions.
\end{itemize}

\subsection{Sub-Microsecond Serving Latency}
To evaluate real-time deployment feasibility, we benchmarked the inference execution speed across multiple hardware backends and execution modes. As reported in Table~\ref{tab:latency_benchmarks}:

\begin{table}[t]
\centering
\footnotesize
\setlength{\tabcolsep}{2.5pt}
\caption{Inference Latency and Throughput Benchmarks.}
\label{tab:latency_benchmarks}
\resizebox{\columnwidth}{!}{%
\begin{tabular}{lccc}
\toprule
\textbf{Engine / Subsystem} & \textbf{Batch} & \textbf{Mean Latency} & \textbf{Throughput} \\
\midrule
\textbf{HDC Holographic Core} & 1 & \textbf{0.293\,$\mu$s} & \textbf{3,407,503 ops/s} \\
eBPF In-Kernel Push & 1 & 0.406\,$\mu$s & 2,461,665 ev/s \\
Native Struct Pack/Unpack & 1 & 0.084\,$\mu$s & 11,850,301 st/s \\
TorchScript JIT (CPU) & 1 & 0.196\,ms & 5,111 seq/s \\
TorchScript JIT (CPU) & 32 & 0.672\,ms & 47,619 seq/s \\
TorchScript JIT (Apple MPS) & 1,024 & 3.905\,ms & \textbf{262,211 seq/s} \\
\bottomrule
\end{tabular}%
}
\end{table}

The HDC Holographic Core achieves an average execution latency of \textbf{0.293 $\mu$s} ($0.29\,\mu\text{s}$), sustaining over \textbf{3.4 million operations per second} per CPU core. Numerical equivalence between native PyTorch execution and the compiled TorchScript JIT graph was strictly verified: $\max |y_{\text{native}} - y_{\text{jit}}| < 4.77 \times 10^{-6}$.

\subsection{Adversarial Chaos Gym Self-Play}
In the distributed Chaos Gym environment, an Actor-Critic \texttt{RedChaosAgent} subjected a 16-node cluster to concurrent thread pool starvation, $25\%$ packet loss, and Byzantine slowloris delays over 1,000 multi-step tournament episodes.

\begin{figure}[t]
\centering
\footnotesize
\setlength{\tabcolsep}{2.5pt}
\resizebox{\columnwidth}{!}{%
\begin{tabular}{lcc}
\toprule
\textbf{Metric} & \textbf{Without Immune Agent} & \textbf{DeepLog-Infinity} \\
\midrule
Mean Blast Radius & 88.4\% of nodes & \textbf{4.0\% of nodes (-95.5\%)} \\
Cluster P99 Latency & 2,840\,ms & \textbf{1,710\,ms (-39.8\%)} \\
SLA Availability & 61.2\% & \textbf{100.0\% (Zero Breach)} \\
CBF Invariant Breaches & 142 events & \textbf{0 events (Provable)} \\
\bottomrule
\end{tabular}%
}
\caption{Adversarial Chaos Gym tournament results over 1,000 episodes.}
\label{fig:chaos_results}
\end{figure}

As shown in Figure~\ref{fig:chaos_results}, the \texttt{BlueImmuneAgent}, constrained by the QP-CBF safety filter, reduced the mean cluster blast radius from 88.4\% down to \textbf{4.0\%}---a \textbf{95.5\% blast radius reduction}. It simultaneously suppressed P99 tail latency by \textbf{39.8\%} while preserving \textbf{100\% SLA availability} without a single barrier invariant violation.

% ------------------------------------------------------------------------------
% 6. DISCUSSION
% ------------------------------------------------------------------------------
\section{Discussion}
\label{sec:discussion}
\textbf{Operational Trade-Offs in Production:} The dual scoring strategy of DeepLog-Infinity addresses the conflicting requirements of cloud operations:
\begin{itemize}[leftmargin=*]
    \item \textit{Tier-1 On-Call Shifts:} SRE teams require exceptionally high precision to avoid alert fatigue. Here, Consecutive Violations ($g=4, c=2$) delivers 92.90\% precision with only 317 false alarms across 570k blocks.
    \item \textit{Security Audits and Forensics:} Security analysts hunting for stealthy multi-stage intrusions prioritize recall. Here, Cumulative NLL at $\tau = 0.40$ delivers 99.99\% recall, catching 16,836 out of 16,838 attacks.
\end{itemize}

\noindent\textbf{eBPF Kernel Safety Constraints:} Deploying neural decision logic into the Linux kernel requires respecting the BPF verifier's strict computational limits ($10^6$ instructions, no unbounded loops). By mapping sequence trajectories to 10,000-bit binary hypervectors, DeepLog-Infinity replaces complex matrix multiplications with unrolled 64-bit bitwise XOR and popcount operations, easily passing Linux BPF verification while operating at $0.29\,\mu\text{s}$ latency.

\noindent\textbf{Zero-Knowledge Verification Overhead:} The Curve25519 Fiat-Shamir Schnorr NIZK proof generates a compact 64-byte cryptographic witness ($R, s$). Verifying a threat hypervector requires only two elliptic curve scalar multiplications ($s B \stackrel{?}{=} R + e V$), requiring less than $45\,\mu\text{s}$ per node. This enables cross-cloud federation over standard gossip protocols without performance degradation.

% ------------------------------------------------------------------------------
% 7. RELATED WORK
% ------------------------------------------------------------------------------
\section{Related Work}
\label{sec:related_work}
\noindent\textbf{Deep Learning for Log Anomaly Detection:} Du et al.~\cite{du2017deeplog} introduced DeepLog, establishing LSTM next-token prediction over parsed log sequences. LogAnomaly~\cite{meng2019loganomaly} incorporated template-level semantic vectors, while LogRobust~\cite{zhang2019logrobust} utilized Bi-LSTM architectures with attention to handle log message instability. Drain~\cite{he2017drain} optimized online log template extraction using fixed-depth parse trees. DeepLog-Infinity advances this lineage by replacing ad-hoc heuristic thresholds with continuous Cumulative NLL, introducing $\langle\text{EOS}\rangle$ boundary modeling, and achieving a new benchmark SOTA F1-score of 72.23\% on 570k blocks.

\noindent\textbf{Kernel Observability and eBPF:} Modern cloud-native observability leverages eBPF for non-invasive tracing (e.g., Cilium, Pixie). However, existing systems primarily focus on raw packet filtering or metrics collection. DeepLog-Infinity is the first architecture to embed Hyperdimensional Computing (HDC)~\cite{kanerva2009hyperdimensional} directly into eBPF tracepoints, enabling autonomous in-kernel sequence classification with zero disk I/O.

\noindent\textbf{Formal Safety and World Models:} Control Barrier Functions (CBFs)~\cite{ames2019control} and Lyapunov contraction methods~\cite{lyapunov1992general} have been extensively deployed in robotics and autonomous flight. Meanwhile, model-based reinforcement learning with Recurrent State Space Models (RSSM)~\cite{hafner2020dreamer} has achieved remarkable predictive control capabilities. DeepLog-Infinity is the first to bridge these two disciplines within cloud computing, guaranteeing mathematical forward invariance and the total prevention of thundering-herd feedback oscillations during autonomous remediation.

% ------------------------------------------------------------------------------
% 8. CONCLUSION
% ------------------------------------------------------------------------------
\section{Conclusion}
\label{sec:conclusion}
We presented DeepLog-Infinity, an autonomous neuro-immune operating platform for hyperscale cloud observability. DeepLog-Infinity addresses the long-standing bottlenecks of log anomaly detection: replacing disk I/O with $0.29\,\mu\text{s}$ in-kernel eBPF hyperdimensional computing; unifying metrics, events, logs, traces, and profiles under physics-informed Jackson queuing constraints; guaranteeing operational safety through Quadratic Programming Control Barrier Functions and Discrete Lyapunov stability; and enabling privacy-preserving cross-tenant threat sharing via Curve25519 Zero-Knowledge proofs. Rigorous empirical evaluation over all 570,204 HDFS production test blocks establishes a new SOTA F1-score of 72.23\%, an 81.2\% reduction in false alerts, and a 95.5\% blast radius reduction under active adversarial chaos, paving the way for provably safe, autonomous cloud infrastructure.

% ------------------------------------------------------------------------------
% BIBLIOGRAPHY
% ------------------------------------------------------------------------------
\bibliographystyle{plain}
\bibliography{references}

\end{document}
